%% =============================================================
%% Section 3: Welfare Comparison
%% =============================================================

\section{Welfare Comparison}
\label{sec:welfare}

We compare the \emph{expected} total loss $\mathbb{E}[\mathcal{L}]$ under forward
guidance and opacity, where expectations are taken over the realisations of both shocks
$(\varepsilon_d, \varepsilon_c)$, and---in the opacity case---over the CB's type.

\subsection{Period-2 Losses}

The period-2 loss for a central bank of type $\lambda$ under either regime is
$\mathcal{L}_2(\lambda) = (y_2^*)^2 + \lambda(\pi_2^*)^2$.
Substituting~\eqref{eq:y2star_clean} and~\eqref{eq:pi2star_clean}:
\begin{equation}
  \mathcal{L}_2(\lambda)
    = \frac{(1 + \lambda^2\kappa^2)\varepsilon_{d,2}^2
            - 2\lambda\kappa\varepsilon_{d,2}\varepsilon_c
            + \lambda(1 + \kappa^2)\varepsilon_c^2}{(1+\lambda\kappa^2)^2}.
  \label{eq:L2_realized}
\end{equation}
Taking expectations over $(\varepsilon_d, \varepsilon_c)$ with
$\mathbb{E}[\varepsilon_d] = \mathbb{E}[\varepsilon_c] = 0$,
$\mathrm{Var}(\varepsilon_d) = \sigma_d^2$, $\mathrm{Var}(\varepsilon_c) = \sigma_c^2$,
and $\mathrm{Cov}(\varepsilon_d,\varepsilon_c) = 0$:
\begin{equation}
  \mathbb{E}[\mathcal{L}_2(\lambda)]
    = \frac{(1+\lambda^2\kappa^2)\sigma_d^2
            + \lambda(1+\kappa^2)\sigma_c^2}{(1+\lambda\kappa^2)^2}.
  \label{eq:EL2_type}
\end{equation}

\subsection{Exchange Rate Misalignment and Period-1 Losses}

Under forward guidance, private agents know the central bank's type and form correct
expectations of the period-2 rate. The period-1 exchange rate is:
\begin{equation}
  q_1^{\mathrm{FG}}(\lambda) = -i_1 - \phi_c(\lambda)\,\varepsilon_c.
  \label{eq:q1_FG}
\end{equation}

Under opacity, agents use the prior-weighted average:
\begin{equation}
  q_1^{\mathrm{OP}} = -i_1 - \bar{\phi}_c\,\varepsilon_c.
  \label{eq:q1_OP}
\end{equation}

Define the \emph{exchange rate misalignment} under opacity as:
\begin{equation}
  \Delta q_1(\lambda) \equiv q_1^{\mathrm{FG}}(\lambda) - q_1^{\mathrm{OP}}
    = -\bigl(\phi_c(\lambda) - \bar{\phi}_c\bigr)\varepsilon_c.
  \label{eq:misalignment}
\end{equation}
Since $\phi_c(\lambda)$ is strictly increasing, $\Delta q_1(\lambda_H) < 0$ and
$\Delta q_1(\lambda_L) > 0$: the exchange rate is too weak (depreciated) under opacity
when the true type is hawkish and too strong (appreciated) when the true type is dovish.

This misalignment feeds into the period-1 IS curve and NKPC (equations
\eqref{eq:IS1_reduced}--\eqref{eq:NKPC1_r}), generating a period-1 welfare loss under
opacity relative to forward guidance.

\subsection{Main Result}

\begin{proposition}[Forward Guidance Weakly Dominates Opacity]
\label{prop:welfare}
Let $\mathcal{L}^r$ denote total expected loss under regime $r \in \{\mathrm{FG},
\mathrm{OP}\}$, where the expectation is taken over shocks and, under opacity, over
the prior distribution of types. Then:
\begin{equation}
  \mathbb{E}[\mathcal{L}^{\mathrm{FG}}] \leq \mathbb{E}[\mathcal{L}^{\mathrm{OP}}],
  \label{eq:welfare_result}
\end{equation}
with strict inequality whenever $p \in (0,1)$ and $\lambda_H \neq \lambda_L$ and
$\sigma_c^2 > 0$.
\end{proposition}

\paragraph{Proof sketch.}
The total expected loss under each regime decomposes as:
\begin{equation}
  \mathbb{E}[\mathcal{L}^r]
    = \mathbb{E}[\mathcal{L}_1^r] + \beta\,\mathbb{E}[\mathcal{L}_2^r].
  \label{eq:total_loss_decomp}
\end{equation}

\noindent\textit{Period-2 comparison.}
Under FG, the period-2 policy is the same function of shocks as under OP (the bank
re-optimises in period 2 regardless of the period-1 announcement). However, under FG
the central bank's known type is $\lambda$, so:
\begin{equation}
  \mathbb{E}[\mathcal{L}_2^{\mathrm{FG}}] = \mathbb{E}[\mathcal{L}_2(\lambda)].
\end{equation}
Under OP, the ex-ante expected period-2 loss integrates over the type distribution:
\begin{equation}
  \mathbb{E}[\mathcal{L}_2^{\mathrm{OP}}]
    = p\,\mathbb{E}[\mathcal{L}_2(\lambda_H)]
    + (1-p)\,\mathbb{E}[\mathcal{L}_2(\lambda_L)].
\end{equation}
Both expressions equal the same weighted average \emph{ex ante}---period-2 losses are
identical in expectation across regimes. The difference arises solely from the
period-1 losses.

\noindent\textit{Period-1 comparison.}
From~\eqref{eq:IS1_reduced} and~\eqref{eq:NKPC1_r}, the period-1 output gap and
inflation under regime $r$ are functions of $\mathbb{E}_1^r[i_2]$. Under FG,
$\mathbb{E}_1^{\mathrm{FG}}[i_2] = \phi_c(\lambda)\varepsilon_c$ (correct). Under OP,
$\mathbb{E}_1^{\mathrm{OP}}[i_2] = \bar{\phi}_c\varepsilon_c$ (an average that may
diverge from the true value). For a given $i_1$, the period-1 loss under opacity
exceeds that under FG by:
\begin{align}
  \mathcal{L}_1^{\mathrm{OP}} - \mathcal{L}_1^{\mathrm{FG}}
  &\propto \bigl(\phi_c(\lambda) - \bar{\phi}_c\bigr)^2 \varepsilon_c^2.
  \label{eq:L1_gap}
\end{align}
Taking expectations over $\varepsilon_c$ and the prior over $\lambda$:
\begin{equation}
  \mathbb{E}[\mathcal{L}_1^{\mathrm{OP}}] - \mathbb{E}[\mathcal{L}_1^{\mathrm{FG}}]
  \propto \mathrm{Var}_\lambda\bigl[\phi_c(\lambda)\bigr]\cdot\sigma_c^2 \geq 0,
  \label{eq:L1_gap_expected}
\end{equation}
where $\mathrm{Var}_\lambda[\phi_c(\lambda)] = p(1-p)(\phi_c(\lambda_H) -
\phi_c(\lambda_L))^2 > 0$ whenever $p \in (0,1)$ and $\lambda_H \neq \lambda_L$.
This is strictly positive when $\sigma_c^2 > 0$, completing the argument. \qed

\paragraph{Economic interpretation.}
The welfare gain from forward guidance has two additive components:

\begin{enumerate}[label=(\roman*)]
  \item \textbf{Exchange rate channel.} By revealing its type, the central bank allows
        markets to correctly price the path of the real exchange rate from period 1
        onwards. Under opacity, the exchange rate is anchored to the prior-weighted
        average response coefficient $\bar{\phi}_c$, generating a misalignment
        $\Delta q_1(\lambda)$ that distorts period-1 allocations. This channel is
        specific to the open-economy setting.

  \item \textbf{Expectation coordination channel.} Correct knowledge of $\phi_c(\lambda)$
        also reduces the variance of expected future inflation $\mathbb{E}_1[\pi_2]$
        and future output $\mathbb{E}_1[y_2]$, tightening the period-1 IS curve and
        NKPC around their optimal values.
\end{enumerate}

Both channels are activated only when cost-push shocks are present ($\sigma_c^2 > 0$)
and genuine type uncertainty exists ($p \in (0,1)$, $\lambda_H \neq \lambda_L$). In
the absence of cost-push shocks, demand shocks are accommodated identically by all
types, $\phi_c$ plays no role, and the two regimes are welfare-equivalent.

\subsection{The Exchange Rate Amplification Effect}
\label{sec:amplification}

A key prediction is that the welfare advantage of forward guidance is \emph{larger}
in more open economies. To see this, note that the misalignment term in
\eqref{eq:L1_gap} involves $\bigl(\phi_c(\lambda) - \bar{\phi}_c\bigr)^2$, where
$\phi_c(\lambda) = \lambda\kappa/[\tilde{\sigma}(1+\lambda\kappa^2)]$.

\begin{corollary}[Openness and the Value of Transparency]
\label{cor:openness}
The welfare gain from forward guidance is decreasing in the degree of openness
$\alpha$ (holding $\sigma$ fixed).
\end{corollary}

\begin{proof}
Since $\tilde{\sigma} = \sigma + \alpha$, increasing $\alpha$ raises $\tilde{\sigma}$,
which reduces $\phi_c(\lambda)$ for all $\lambda$. As $\tilde{\sigma} \to \infty$, all
$\phi_c(\lambda) \to 0$ and the welfare gap~\eqref{eq:L1_gap_expected} vanishes.
The amplification is therefore \emph{non-monotone}: openness amplifies the
transmission of the policy rate to output (via $\tilde{\sigma}$), but it also
compresses the dispersion of reaction functions across types. \qed
\end{proof}

\noindent\textbf{Remark.} Corollary~\ref{cor:openness} has a substantive implication:
very open economies may have less to gain from formal forward guidance, because the
exchange rate channel already transmits the policy stance efficiently---the market can
``read'' the central bank's intentions from the current exchange rate. This is
consistent with the observation that several highly open economies (e.g.\ Hong Kong,
Singapore) operate under monetary frameworks that rely less on explicit rate-path
guidance.
